\section[Kiểm chuẩn]{Kiểm chuẩn (Benchmark)}

Trong phần này, các mô hình được đánh giá và so sánh dựa trên nhiều tiêu chí
khác nhau như thời gian huấn luyện, mức sử dụng bộ nhớ trong quá trình huấn
luyện, cũng như thời gian và bộ nhớ khi suy luận. Mục tiêu của việc kiểm chuẩn
là nhằm đánh giá hiệu năng tính toán của từng mô hình trong các điều kiện thực
nghiệm giống nhau.

\subsection{Hồi quy}


Kết quả kiểm chuẩn của các mô hình hồi quy được trình bày trong
\tablename~\ref{tab:regression-benchmark-part1} và
\tablename~\ref{tab:regression-benchmark-part2}. Kết quả được tính trên 5 lần
chạy. Các tiêu chí đánh giá gồm trung bình và độ lệch chuẩn thời gian (theo đơn
vị giây) và bộ nhớ (tính theo đơn vị byte) quá trình huấn luyện và suy luận.

\begin{table}[htbp]
  \centering
  \caption{Kết quả kiểm chuẩn các mô hình hồi quy (phần 1)}
  \label{tab:regression-benchmark-part1}
  \import{\tablespath}{regression-benchmark-part1}
\end{table}

\begin{table}[htbp]
  \centering
  \caption{Kết quả kiểm chuẩn các mô hình hồi quy (phần 2)}
  \label{tab:regression-benchmark-part2}
  \import{\tablespath}{regression-benchmark-part2}
\end{table}

Từ các kết quả kiểm chuẩn, có thể thấy sự khác biệt rõ rệt giữa các mô hình về
thời gian huấn luyện và mức sử dụng bộ nhớ.

Các mô hình tuyến tính đơn giản như OLS và WLS có thời gian huấn luyện rất
thấp, gần như tức thời, đồng thời mức tiêu thụ bộ nhớ cũng nhỏ. Điều này cho
thấy đây là các mô hình có độ phức tạp tính toán thấp, phù hợp với các bài toán
cần tốc độ xử lý cao.

Ngược lại, MBGD có thời gian huấn luyện lớn hơn đáng kể so với OLS và WLS do
phải cập nhật tham số theo từng lô dữ liệu, kéo theo chi phí tính toán tăng
lên. Tuy nhiên, mức bộ nhớ sử dụng vẫn tương đối ổn định.

Các mô hình Ridge, Lasso và Elastic Net có thời gian huấn luyện cao hơn OLS
nhưng vẫn ở mức chấp nhận được. Sự khác biệt giữa ba mô hình này không quá lớn,
tuy nhiên Elastic Net có xu hướng tốn thời gian hơn nhẹ do kết hợp đồng thời
hai dạng chính quy.

Về suy luận, tất cả các mô hình hồi quy đều có thời gian suy luận rất nhỏ, gần
như không đáng kể, cho thấy khả năng dự đoán nhanh sau khi huấn luyện.

\subsection{Phân lớp}

Kết quả kiểm chuẩn của các mô hình phân lớp được trình bày trong
\tablename~\ref{tab:classification-benchmark-part1} và
\tablename~\ref{tab:classification-benchmark-part2}.

\begin{table}[htbp]
  \centering
  \caption{Kết quả kiểm chuẩn các mô hình phân lớp (phần 1)}
  \label{tab:classification-benchmark-part1}
  \import{\tablespath}{classification-benchmark-part1}
\end{table}

\begin{table}[htbp]
  \centering
  \caption{Kết quả kiểm chuẩn các mô hình phân lớp (phần 2)}
  \label{tab:classification-benchmark-part2}
  \import{\tablespath}{classification-benchmark-part2}
\end{table}

Trong nhóm mô hình phân lớp, có sự khác biệt rõ rệt hơn về cả thời gian huấn
luyện lẫn bộ nhớ.

LDA và QDA có thời gian huấn luyện tương đối nhỏ, trong đó QDA nhanh hơn LDA
nhưng lại có thời gian suy luận lớn hơn đáng kể. Điều này phản ánh sự đánh đổi
giữa độ phức tạp mô hình và chi phí dự đoán.

Softmax và Weighted Softmax có thời gian huấn luyện cao hơn so với LDA và QDA,
trong đó Weighted Softmax là mô hình tốn thời gian huấn luyện nhiều nhất. Tuy
nhiên, chi phí suy luận của các mô hình này khá thấp và ổn định.

Về bộ nhớ, LDA có mức tiêu thụ bộ nhớ cao hơn so với các mô hình còn lại, trong
khi Softmax và Weighted Softmax sử dụng bộ nhớ ở mức trung bình nhưng ổn định
giữa các lần chạy.
